\section{Introduction}\label{sec:introduction}

A cell-fate model may predict apoptosis during sustained stimulation without
specifying whether that response still depends on the stimulus. This distinction
matters when two mechanistic hypotheses predict the same ordinary endpoints:
the informative comparison may concern the futures remaining after an
intervention, rather than their original list of attractors. Death-receptor
signaling provides a documented setting. Calzone et al. compared a model
containing CASP3-to-CASP8 positive feedback with a variant lacking that
interaction, and analyzed the persistence of apoptosis after transient versus
sustained ligand exposure.\cite{Calzone2010} We use that published phenomenon
to ask a narrower, executable question: \emph{after a shared history and with
the same readouts, which future responses can each variant still reproduce
when TNF is withdrawn?}

Qualitative biological models make such hypotheses executable without requiring
fully parameterized kinetic descriptions.\cite{ChaouiyaPetriBio2007,Li2006,Fisher2007}
Their comparison nevertheless depends on initial conditions, update semantics,
and observable events. Structural similarity characterizes wiring, not the
availability of actions at a reachable state.\cite{Szawulak2022} Trace comparison
adds event order but can discard where alternatives cease to coexist.
Example~\ref{ex:branching} makes that mathematical distinction explicit with
early and late commitment to a choice. Here it motivates the biological
search for common observations followed by different intervention responses;
it does not predetermine that a biological pair must have equal global traces.

We apply established weak simulation and bisimulation\cite{Milner1989,Sangiorgi2011}
to finite executable models. In $X\weakpre Y$, the right-hand model must match
each step of the left-hand model and retain that ability at the resulting
state pair. Bisimulation uses one consistent relation in both directions.
The observational interface is fixed before these calculations, and hidden
internal events may be matched silently. The comparison map is
\begin{equation}\label{eq:comparison-map}
 (N_1,N_2,\Sigma)\longmapsto(\LTS(N_1),\LTS(N_2))
 \longmapsto C_\Sigma(\LTS(N_1),\LTS(N_2)),
\end{equation}
where $C_\Sigma$ reports the supported relation without assigning a biological
quality ranking to its classes.

The primary case compares the original death-receptor model, denoted DR-FB+,
with the documented feedback-deletion variant, DR-FB-. Under sustained TNF,
both have the same three reachable stable fates. We identify a shared
CASP3-positive state with exactly matching sustained continuations but different
responses to withdrawal. The full withdrawal-aware systems also have explicit
trace counterexamples in both directions. Thus the demonstrated value is an
intervention-sensitive characterization of commitment in endpoint-matched
models, not exclusive superiority over trace analysis once the discriminating
intervention is included. A verified finite counterexample also refutes global
sustained trace equivalence, so trace analysis can distinguish the complete
baseline systems as well. Only the reverse sustained trace inclusion remains
undecided; no bounded trace distance is substituted for an exact decision.

A complementary GIM (genome-instability/DNA-repair) example concerns a different
question: how correspondence changes when the interface itself changes.
Animal/human and \tas{} encodings share broad damage-response functions but
contain independently documented terminal differences.\cite{JacksonBartek2009,
BlackfordJackson2017,Yoshiyama2013,DeSchutter2007,Yi2014,Adachi2011}
Their structural similarity remains 0.9976 while correspondence is lost at the
first of three declared interfaces that exposes those differences. GIM is
retained as an illustration of interface-conditioned correspondence, not the
principal biological evidence or a search for a unique biological boundary.

The contribution is a reproducible connection between an explicit biological
perturbation question, reachable future sets, and checked relational witnesses.
It recovers an already published feedback-dependent persistence phenomenon;
neither the mechanism nor the idea of ligand withdrawal is a new discovery.
Supporting curated cases and HPN-DREAM/CASPOTS response data retain their
separate evidential roles.\cite{Hill2016,Razzaq2018}
